%pdf-latex

%header nur vom MÜSLI-Vortrag geklaut
% !TEX root = vortrag.tex
% !TEX encoding = UTF-8 Unicode

\documentclass[t, ngerman]{beamer} %compress,

%% Pakete laden...
  \usepackage[T1]{fontenc}
  \usepackage[utf8]{inputenc}
  \usepackage{
      babel,
      bookmark,
      booktabs,
%      blindtext,
      colortbl,
%      eurosym,
      graphicx,
	  hyperref,
%      libertine,
      microtype,
      pifont,
      pgfpages,
      tikz,
%      xspace,
  }


%% Design festlegen...
  \mode<presentation>{
%      \useoutertheme[subsection=false]{smoothbars}
      \useinnertheme{rectangles} % rectangles, circles, rounded
      \usecolortheme[RGB={153,0,0}]{structure}
      \definecolor{unihd}{RGB}{153,0,0}
      \definecolor{dark}{RGB}{115,0,0}
      \definecolor{light}{RGB}{241,229,229}
      \usecolortheme{whale}
		 \usecolortheme{orchid}
%	   \usecolortheme{beaver}
%      \setbeamercovered{transparent}
      \beamertemplatenavigationsymbolsempty
%      \setbeameroption{show notes on second screen}
      \setbeamertemplate{note page}[plain]
      \logo{\includegraphics[width=3.5cm]{fs-logo-small}}

  }

%% nützliche Definitionen...
  \graphicspath{{media/}}

%% Titelinformationen...
  \title[Studienverwaltung$\mu$]{Das Müsli und das LSF\\\small oder: wer wie was wo Stundenplan}
  \author[
	koebi
  ]{
	Jakob Schnell\\{\scriptsize\url{koebi@mathphys.stura.uni-heidelberg.de}}
  }
%  \institute{  \includegraphics[width=5.5cm]{fs-logo-big} }

  \date{\vspace*{-2em}\\ 01. Oktober 2018}

  \hypersetup{
      pdfauthor={Jakob Schnell},
      pdftitle={Müsli-Vortrag},
      pdfsubject={hihi},
      pdfkeywords={1},
      pdfpagelayout={SinglePage},
  }
  


\newenvironment{rcases}{%
  \left.\renewcommand*\lbrace.%
  \begin{cases}}%
{\end{cases}\right\rbrace}


\begin{document}

\begin{frame}[fragile]
    \maketitle{}
\end{frame}

\begin{frame}{Inhalt des Vortrags}
    \begin{minipage}[t]{0.515\textwidth}
        \tableofcontents[hideallsubsections, sections={1-5}]
    \end{minipage}
    \begin{minipage}[t]{0.475\textwidth}
        \tableofcontents[hideallsubsections, sections={6-11}]
    \end{minipage}
\end{frame}

\begin{frame}{Wer will mir hier was erzählen!?}
    \vfill
    \begin{center}
        {\Large Janina Rastetter} \\
        \mail{jrastetter@mathi.uni-heidelberg.de} \\
        \vspace{1em}
        {\Large Christian Heusel } \\
        \mail{chris@mathphys.stura.uni-heidelberg.de} \\
    \end{center}
\end{frame}

\begin{frame}{Präambel}
    \Large
    \begin{center}
        \only<1->{%
            {\huge \DejaSans{} 😇} Es wird die Folien zum Download geben! {\huge \DejaSans{} 😇}\\[0.7em]
        }
        \only<1>{\includegraphics[scale=0.6]{images/xkcd_1.png}}
        \only<2->{%
            {\huge \DejaSans{} 😺} Stellt Fragen, im Währenden oder am Ende! {\huge \DejaSans{} 😺}\\[0.7em]
        }
        \only<3->{%
            {\huge \DejaSans{} 😱} Es wird viele Informationen geben! {\huge \DejaSans{} 😱}\\[0.7em]
        }
        \only<4->{%
            \includegraphics[]{images/lets_go.png}
        }
    \end{center}
\end{frame}

%%%%%%%%%%%%%%%%%%%%%%%%%%%%%%%%%%%%%%%%%%%%%%%%%%%%%%%%%%%%
% Digitale Infrastruktur
%%%%%%%%%%%%%%%%%%%%%%%%%%%%%%%%%%%%%%%%%%%%%%%%%%%%%%%%%%%%

\section{Digitale Infrastruktur}
\begin{frame}{Digitale Infrastruktur}
    \large
    \begin{itemize}
        \item{WLAN}
        \item{Drucken}
        \item{VPN}
        \item{CIP-Pool}
        \item{Uni-Mail-Adresse}
    \end{itemize}
\end{frame}

%%%%%%%%%%%%%%%%%%%%%%%%%%%%%%%%%%%%%%%%%%%%%%%%%%%%%%%%%%%%
% WLAN
%%%%%%%%%%%%%%%%%%%%%%%%%%%%%%%%%%%%%%%%%%%%%%%%%%%%%%%%%%%%

\subsection{WLAN}
\begin{frame}{Eduroam}
\end{frame}

\begin{frame}{Eduroam einrichten -- leicht gemacht}
\end{frame}

%%%%%%%%%%%%%%%%%%%%%%%%%%%%%%%%%%%%%%%%%%%%%%%%%%%%%%%%%%%%
% Drucken
%%%%%%%%%%%%%%%%%%%%%%%%%%%%%%%%%%%%%%%%%%%%%%%%%%%%%%%%%%%%

\subsection{Drucken}
\begin{frame}{Drucken}
\end{frame}

%%%%%%%%%%%%%%%%%%%%%%%%%%%%%%%%%%%%%%%%%%%%%%%%%%%%%%%%%%%%
% VPN
%%%%%%%%%%%%%%%%%%%%%%%%%%%%%%%%%%%%%%%%%%%%%%%%%%%%%%%%%%%%

\subsection{VPN}
\begin{frame}{Was ist denn VPN?}
    \large
    \begin{itemize}
        \item \textbf{V}irtual \textbf{P}rivate \textbf{N}etwork
        \item Ein VPN ermöglicht den (Fern-)Zugriff auf andere Netze
        \item Alle Daten werden für den Transport verschlüsselt
        \item hat erstmal nichts mit einem VPN Anbieter zu tun \\
              (NordVPN, TunnelBear, ExpressVPN)
        \item bla bla bla
    \end{itemize}
\end{frame}

\begin{frame}{Wozu nützt mir ein VPN?}
\end{frame}

\begin{frame}{Wie sieht das aus?}
    \begin{center}
        \only<1>{%
            \vspace*{-1.1cm}
            \includegraphics[scale=2.3]{images/vpn_1.png}
        }
        \only<2>{%
            \vspace*{-0.7cm}
            \includegraphics[scale=2.3]{images/vpn.png}
        }
    \end{center}
\end{frame}

%%%%%%%%%%%%%%%%%%%%%%%%%%%%%%%%%%%%%%%%%%%%%%%%%%%%%%%%%%%%
% CIP-Pool
%%%%%%%%%%%%%%%%%%%%%%%%%%%%%%%%%%%%%%%%%%%%%%%%%%%%%%%%%%%%

\subsection{CIP-Pool}
\begin{frame}{CIP-Pool}
\end{frame}

%%%%%%%%%%%%%%%%%%%%%%%%%%%%%%%%%%%%%%%%%%%%%%%%%%%%%%%%%%%%
% Uni-Mail-Adresse
%%%%%%%%%%%%%%%%%%%%%%%%%%%%%%%%%%%%%%%%%%%%%%%%%%%%%%%%%%%%

\subsection{Uni-Mail-Adresse}
\begin{frame}{Uni-Mail-Adresse}
\end{frame}

%%%%%%%%%%%%%%%%%%%%%%%%%%%%%%%%%%%%%%%%%%%%%%%%%%%%%%%%%%%%
% Wichtige Websites
%%%%%%%%%%%%%%%%%%%%%%%%%%%%%%%%%%%%%%%%%%%%%%%%%%%%%%%%%%%%

\section{Wichtige Websites}
\begin{frame}{Wichtige Websites}
    \large Die wichtigsten Websites im Unikontext sind:
    \begin{itemize}
        \item{LSF}
        \item{MÜSLI}
        \item{Moodle}
        \item{MaMpf}
    \end{itemize}
\end{frame}

%%%%%%%%%%%%%%%%%%%%%%%%%%%%%%%%%%%%%%%%%%%%%%%%%%%%%%%%%%%%
% LSF
%%%%%%%%%%%%%%%%%%%%%%%%%%%%%%%%%%%%%%%%%%%%%%%%%%%%%%%%%%%%

\subsection{LSF}
\begin{frame}{Das LSF -- Lehre, Studium und Forschung}

    \large \url{https://lsf.uni-heidelberg.de} \\
    \begin{minipage}[t]{0.515\textwidth}
        \begin{itemize}
            \item{Veranstaltungssuche}
            \item{Studienbescheinigung}
            \item{Rückmeldung}
            \item{Bafög-Daten}
            \item{Noten}
            \item{Stundenplan}
        \end{itemize}
    \end{minipage}
    \begin{minipage}[t]{0.475\textwidth}
        \vspace{0.1cm}
        \begin{center}
            \qrcode[height=0.65\textwidth, link]{https://lsf.uni-heidelberg.de}
        \end{center}
    \end{minipage}
\end{frame}

\begin{frame}{Veranstaltungssuche}
    \only<1>{
          Welche Veranstaltungen ihr belegen müsst/könnt, entnehmt ihr dem Modulhandbuch zu eurem Studiengang. Dazu wird es einen gesonderten Vortrag geben.\\
          \vspace{0.5cm}
          Achtung!\\
          \begin{itemize}
                \item{Nach (Pro-)Seminaren haltet ihr am besten in den Aushängen im Mathematikon oder im MÜSLI am Ende der Vorlesungszeit Ausschau (Vorbesprechungen finden teilweise in der vorletzten oder letzten Vorlesungswoche statt)}
                \item{Über Übungsgruppen informiert ihr euch im MÜSLI}
                \item{Für alles andere gibt es das LSF}
          \end{itemize}
    }
    \only<2>{
          \begin{figure}
                \centering
                \includegraphics[scale=0.3]{images/lsf01.jpg}
          \end{figure}
          Achtet darauf, dass das richtige Semester eingestellt ist, bevor ihr nach Veranstaltungen sucht.
    }
    \only<3>{
          \begin{figure}
                \centering
                \includegraphics[scale=0.25]{images/lsf02.jpg}
          \end{figure}
    }
    \only<4>{
          Zwei Möglichkeiten, im LSF nach Veranstaltungen zu suchen
          \begin{enumerate}
                \item{Veranstaltungsangebot einer Fakultät durchsuchen (Was wird angeboten?)}
                \item{Nach konkreter Veranstaltung über die Suchfunktion suchen (Wird eine bestimmte Veranstaltung angeboten? Wann und wo findet sie statt?)}
          \end{enumerate}
    }
\end{frame}

\begin{frame}{Veranstaltungssuche über Fakultät -- Fach -- Studiengang}
    \only<1>{%}
        \begin{figure}
            \centering
            \includegraphics[scale=0.3]{images/lsf03.jpg}
        \end{figure}
    }
    \only<2>{%}
        \begin{figure}
            \centering
            \includegraphics[scale=0.25]{images/lsf04.jpg}
        \end{figure}
    }
    \only<3>{%}
        \begin{figure}
            \centering
            \includegraphics[scale=0.25]{images/lsf05.jpg}
        \end{figure}
    }
    \only<4>{%}
        \begin{figure}
            \centering
            \includegraphics[scale=0.25]{images/lsf06.jpg}
        \end{figure}
    }
    \only<5>{%}
        \begin{figure}
            \centering
            \includegraphics[scale=0.25]{images/lsf07.jpg}
        \end{figure}
    }
    \only<6>{%}
        \begin{figure}
            \centering
            \includegraphics[scale=0.25]{images/lsf08.jpg}
        \end{figure}
    }
\end{frame}

\begin{frame}{Veranstaltungssuche über die Suchfunktion}
    \only<1>{%
        \begin{figure}
            \centering
            \includegraphics[scale=0.25]{images/lsf09.jpg}
        \end{figure}
    }
    \only<2>{%
        \begin{figure}
           \centering
           \includegraphics[scale=0.3]{images/lsf10.jpg}
        \end{figure}
    }
\end{frame}

\begin{frame}{Studienbescheinigung, Bafög und Rückmeldung}
    Für diese Funktionen müsst ihr euch anmelden.
    \begin{figure}
        \centering
        \includegraphics[scale=0.3]{images/lsf12.jpg}
    \end{figure}
\end{frame}

\begin{frame}{Rückmeldung}
    \begin{itemize}
        \item{Erforderlich, wenn ihr weiter studieren wollt}
        \item{Rückmeldung = Semesterbeiträge bezahlen}
        \item{Aufforderung dazu: über die Uni-Mail-Adresse jeweils Mitte Januar und Mitte Juni}
    \end{itemize}
\end{frame}

\begin{frame}{Noten}
    \begin{figure}
        \centering
        \includegraphics[scale=0.3]{images/lsf13.jpg}
    \end{figure}
    Eure Noten werden -- mit etwas Verzögerung -- vom jeweiligen fachinternen System hierher übertragen. Auch hierfür müsst ihr euch anmelden.
\end{frame}

\begin{frame}{Stundenplan}
    \begin{figure}
        \centering
        \includegraphics[scale=0.3]{images/lsf14.jpg}
    \end{figure}
    Achtung! Auf diese Weise meldet ihr euch NICHT für Kurse an.
\end{frame}

\begin{frame}{Studenplan}
    \begin{figure}
        \centering
        \includegraphics[scale=0.25]{images/lsf15.jpg}
    \end{figure}
\end{frame}

%%%%%%%%%%%%%%%%%%%%%%%%%%%%%%%%%%%%%%%%%%%%%%%%%%%%%%%%%%%%
% MÜSLI
%%%%%%%%%%%%%%%%%%%%%%%%%%%%%%%%%%%%%%%%%%%%%%%%%%%%%%%%%%%%

\subsection{MÜSLI}
\begin{frame}{MÜSLI -- \normalsize Mathematisches Übungsgruppen- und Scheinlisten-Interface}

    \large \url{https://muesli.mathi.uni-heidelberg.de/}

    \begin{minipage}[t]{0.7\textwidth}

    \begin{itemize}
        \item Eintragung in Übungsgruppen
        \item Einsehen von Zettelpunkten
        \item E-Mail-Adressen der Tutor*innen
        \item Klausuranmeldung
        \item Noten
    \end{itemize}
    \end{minipage}
    \begin{minipage}[t]{0.28\textwidth}
        \vspace*{0em}
        \begin{center}
            \qrcode[height=0.9\textwidth]{https://muesli.mathi.uni-heidelberg.de/}
        \end{center}
    \end{minipage}
\end{frame}

\begin{frame}{MÜSLI -- \normalsize Mathematisches Übungsgruppen- und Scheinlisten-Interface}
    \Large {\LARGE \DejaSans{} ⚠} Achtung {\LARGE \DejaSans{}⚠} \\
    \normalsize
    \begin{enumerate}
        \item{Ihr müsst euch mit eurer Uni-Mail-Adresse registrieren.}
        \item{In anderen Fächern gibt es andere Anmeldemodalitäten. Wenn ihr Kurse in einem anderen Fach belegen wollt/müsst, informiert euch frühzeitig(!) darüber, wie ihr euch in diesem Fach zu Veranstaltungen anmelden könnt.}
    \end{enumerate}
\end{frame}

%%%%%%%%%%%%%%%%%%%%%%%%%%%%%%%%%%%%%%%%%%%%%%%%%%%%%%%%%%%%
% Moodle
%%%%%%%%%%%%%%%%%%%%%%%%%%%%%%%%%%%%%%%%%%%%%%%%%%%%%%%%%%%%

\subsection{Moodle}
\begin{frame}{Moodle -- uniweite E-Learningplattform}

    Link: \url{https://elearning2.uni-heidelberg.de/}

    \begin{center}
        \qrcode{https://elearning2.uni-heidelberg.de/}
    \end{center}

    \begin{itemize}
        \item Anmeldung über Uni-ID
        \item{Kurse, in die ihr euch mit einem Einschreibeschlüssel eintragen könnt; Dozent*innen geben den Einschreibeschlüssel i.d.R. in der ersten Sitzung bekannt}
        \item{Materialien zu Veranstaltungen wie Übungsblätter oder Präsentationen}
        \item{Abgaben}
    \end{itemize}

\end{frame}

\begin{frame}{Moodle -- uniweite E-Learningplattform}
    \only<1>{%}
        \begin{figure}
            \centering
            \includegraphics[scale=0.3]{images/moodle01.jpg}
        \end{figure}
    }
    \only<2>{%}
        \begin{figure}
            \centering
            \includegraphics[scale=0.3]{images/moodle02.jpg}
        \end{figure}
    }
    \only<3>{%}
        \begin{figure}
            \centering
            \includegraphics[scale=0.3]{images/moodle03.jpg}
        \end{figure}
    }
    \only<4>{%}
        \begin{figure}
            \centering
            \includegraphics[scale=0.3]{images/moodle04.jpg}
        \end{figure}
    }
    \only<5>{%}
        \begin{figure}
            \centering
            \includegraphics[scale=0.3]{images/moodle05.jpg}
        \end{figure}
    }
\end{frame}

%%%%%%%%%%%%%%%%%%%%%%%%%%%%%%%%%%%%%%%%%%%%%%%%%%%%%%%%%%%%
% MaMpf
%%%%%%%%%%%%%%%%%%%%%%%%%%%%%%%%%%%%%%%%%%%%%%%%%%%%%%%%%%%%

\subsection{MaMpf}
\begin{frame}{\includegraphics[scale=0.072]{images/mampf.png} MaMpf -- Mathematische Medienplattform}

    Link: \url{https://mampf.mathi.uni-heidelberg.de/}

    \begin{center}
        \qrcode{https://mampf.mathi.uni-heidelberg.de/}
    \end{center}

    \begin{itemize}
        \item Speziell fürs Mathelernen entwickelte E-Learning-Plattform
        \item Viele Lernangebote (Vorlesungsvideos, Beispielsvideos, Quizzes, ...)
        \item Werdet ihr benutzen, wenn ihr die LA 1 hört
        \item Richtig gut
    \end{itemize}

\end{frame}

\begin{frame}{Die wichtigsten Websites -- Zusammenfassung}
    \begin{itemize}
        \item LSF: allgemeine Studiumsverwaltung
        \item MÜSLI: Studiumsverwaltung in Mathe und Info
        \item Moodle: uniweite E-Learning-Plattform
        \item MaMpf: E-Learning-Plattform für Mathe
    \end{itemize}
\end{frame}

%%%%%%%%%%%%%%%%%%%%%%%%%%%%%%%%%%%%%%%%%%%%%%%%%%%%%%%%%%%%
% Software/Anleitungen
%%%%%%%%%%%%%%%%%%%%%%%%%%%%%%%%%%%%%%%%%%%%%%%%%%%%%%%%%%%%

\section{Software \& Anleitungen}
\begin{frame}{Software/Anleitungen}
\end{frame}

%%%%%%%%%%%%%%%%%%%%%%%%%%%%%%%%%%%%%%%%%%%%%%%%%%%%%%%%%%%%
% Universitätsbibliothek
%%%%%%%%%%%%%%%%%%%%%%%%%%%%%%%%%%%%%%%%%%%%%%%%%%%%%%%%%%%%

\section{Unibibliothek}
\begin{frame}{Unibibliothek}
    \only<1>{%
        Welche Bibliotheken gibt es an der Uni?
        \begin{itemize}
            \item Institutsbibliotheken
                \begin{itemize}
                    \item Fachspezifisch
                    \item Normalerweise nur Präsenzbibliotheken, d.h. keine Ausleihe möglich
                    \item Mathe-Info-Bib befindet sich im EG des Mathematikons
                \end{itemize}
            \item Zentralbibliotheken
                \begin{itemize}
                    \item  Nicht fachspezifisch
                    \item Hauptbibliothek in der Altstadt (Plöck 107-109, Ausleihe im EG)
                    \item Zweigstelle im Neuenheimer Feld (Im Neuenheimer Feld 368, Ausleihe im 3.OG)
                \end{itemize}
        \end{itemize}
    }
    \only<2>{%
        Beachtet
        \begin{itemize}
            \item {Bibliotheksangebot nutzen}
                \begin{itemize}
                    \item Studierendenausweis = Nutzerausweis
                    \item{Um alle Angebote nutzen zu können, müsst ihr euren Ausweis freischalten \url{https://www.ub.uni-heidelberg.de/service/anmeldung.html}}
                \end{itemize}
            \item Schließfächer
                \begin{itemize}
                    \item Zentralbibliotheken: Ihr braucht eine Zwei-Euro-Münze
                    \item Mathe-Info-Bib: Es gibt Schließfächer mit Zahlencodes oder Schlüsseln; Schließfachschüssel erhaltet ihr bei der Bibliotheksaufsicht gegen ein Pfand (z.B. Ausweis)
                \end{itemize}
        \end{itemize}
    }
    \only<3>{%
        Was bieten euch Bibliotheken?
        \begin{itemize}
            \item{Lernort: allein oder zusammen (Gruppenarbeitsräume reservieren: \url{https://www.ub.uni-heidelberg.de/service/gruppenarbeitsraeume.html})}
            \item{Literatur}
                  \begin{itemize}
                        \item{Finden, nutzen, ausleihen und downloaden}
                        \item{Mehr Literatur über Datenbanken finden}
                        \item{Fernleihe}
                        \item{Teilweise Zugriff auf Online-Angebot von Zeitschriften}
                  \end{itemize}
            \item{Weitere Informationen (z.B. Führungen und Kurse) \url{https://www.ub.uni-heidelberg.de/schulung/Welcome.html}}
        \end{itemize}
    }
\end{frame}

\begin{frame}{Unibiliothek -- HEIDI}
      \only<1>{%}
          Bibliothekskatolog HEIDI
          \begin{itemize}
                \item \url{https://katalog.ub.uni-heidelberg.de/cgi-bin/search.cgi?zweig=}
                \item Medien finden und ihren Standort in Erfahrung bringen oder downloaden
          \end{itemize}
      }
      \only<2>{
            \begin{figure}
                \centering
                \includegraphics[scale=0.35]{images/ub01.jpg}
            \end{figure}
      }
      \only<3>{
            \begin{figure}
                \centering
                \includegraphics[scale=0.2]{images/ub02.jpg}
            \end{figure}
      }
      \only<4>{
            \begin{figure}
                \centering
                \includegraphics[scale=0.2]{images/ub03.jpg}
            \end{figure}
      }
      \only<5>{
            \begin{figure}
                \centering
                \includegraphics[scale=0.2]{images/ub04.jpg}
            \end{figure}
      }
\end{frame}

%%%%%%%%%%%%%%%%%%%%%%%%%%%%%%%%%%%%%%%%%%%%%%%%%%%%%%%%%%%%
% Fachschaftsservices
%%%%%%%%%%%%%%%%%%%%%%%%%%%%%%%%%%%%%%%%%%%%%%%%%%%%%%%%%%%%

\section{Fachschaftsservices}
\begin{frame}{Fachschaftsservices}
\end{frame}

\begin{frame}{Wer will mir hier was erzählen!?}
    \vfill
    \begin{center}
        {\Large Janina Rastetter} \\
        \mail{jrastetter@mathi.uni-heidelberg.de} \\
        \vspace{1em}
        {\Large Christian Heusel } \\
        \mail{chris@mathphys.stura.uni-heidelberg.de} \\
    \end{center}
\end{frame}

\begin{frame}{Vielen Dank fürs Zuhören!}
    \vfill
    \begin{center}
        \Huge Fragen \\[3.5pt]
        \normalsize Serviceangebote Studium
    \end{center}
\end{frame}

\end{document}
